\documentclass[11pt,a4paper]{article}

% ── Packages ─────────────────────────────────────────────────────────────
\usepackage[margin=1in]{geometry}
\usepackage[T1]{fontenc}
\usepackage[utf8]{inputenc}
\usepackage{parskip}
\usepackage{xcolor}
\usepackage{graphicx}
\usepackage{booktabs}
\usepackage{listings}
\usepackage{tcolorbox}
\usepackage{titlesec}
\usepackage{enumitem}
\usepackage{fancyhdr}
\usepackage{caption}
\usepackage{tabularx}
\usepackage[hidelinks,colorlinks=false,pdfborder={0 0 0}]{hyperref}

% ── Colors ───────────────────────────────────────────────────────────────
\definecolor{slappink}{HTML}{E11D74}
\definecolor{softpink}{HTML}{FFF1F5}
\definecolor{codebg}{HTML}{F6F6F8}
\definecolor{decisionblue}{HTML}{1E40AF}
\definecolor{softblue}{HTML}{EFF6FF}
\definecolor{warnamber}{HTML}{B45309}
\definecolor{softamber}{HTML}{FEF3C7}
\definecolor{muted}{HTML}{57534E}

% ── Code listing style ───────────────────────────────────────────────────
\lstdefinestyle{ourstyle}{
  backgroundcolor=\color{codebg},
  basicstyle=\ttfamily\footnotesize,
  breaklines=true,
  breakatwhitespace=true,
  captionpos=b,
  frame=none,
  keywordstyle=\color{decisionblue}\bfseries,
  commentstyle=\color{muted}\itshape,
  stringstyle=\color{slappink},
  showstringspaces=false,
  tabsize=2,
  xleftmargin=0.5em,
  aboveskip=0.6em,
  belowskip=0.6em,
}
\lstset{style=ourstyle}

% ── Custom boxes ─────────────────────────────────────────────────────────
\newtcolorbox{decisionbox}[1][]{
  colback=softblue,
  colframe=decisionblue,
  boxrule=0.6pt,
  arc=1.5pt,
  left=8pt,right=8pt,top=6pt,bottom=6pt,
  title=\textbf{Decision:} #1,
  fonttitle=\small,
  coltitle=white,
  colbacktitle=decisionblue,
}

\newtcolorbox{gotchabox}[1][]{
  colback=softamber,
  colframe=warnamber,
  boxrule=0.6pt,
  arc=1.5pt,
  left=8pt,right=8pt,top=6pt,bottom=6pt,
  title=\textbf{Gotcha:} #1,
  fonttitle=\small,
  coltitle=white,
  colbacktitle=warnamber,
}

\newtcolorbox{screenshotbox}[1][]{
  colback=softpink,
  colframe=slappink,
  boxrule=0.6pt,
  arc=1.5pt,
  left=8pt,right=8pt,top=6pt,bottom=6pt,
  title=\textbf{Screenshot placeholder:} #1,
  fonttitle=\small,
  coltitle=white,
  colbacktitle=slappink,
}

% ── Headers / footers ────────────────────────────────────────────────────
\pagestyle{fancy}
\fancyhf{}
\rhead{\small SLAP Bug Triage --- Final Report}
\lhead{\small Angaraju V}
\cfoot{\thepage}
\renewcommand{\headrulewidth}{0.3pt}

% ── Section formatting ───────────────────────────────────────────────────
\titleformat{\section}{\LARGE\bfseries\color{slappink}}{\thesection}{0.7em}{}
\titleformat{\subsection}{\large\bfseries\color{decisionblue}}{\thesubsection}{0.6em}{}
\titleformat{\subsubsection}{\normalsize\bfseries}{\thesubsubsection}{0.5em}{}

% ── Metadata ─────────────────────────────────────────────────────────────
\hypersetup{
  pdftitle={SLAP Bug Triage --- Final Report},
  pdfauthor={Angaraju V}
}

\begin{document}

% ────────────────────────── TITLE PAGE ──────────────────────────
\begin{titlepage}
\centering
\vspace*{2cm}
{\Huge\bfseries\color{slappink} SLAP Bug Triage\par}
\vspace{0.3cm}
{\Large An agentic bug-triage system for Flipkart's GenAI shopping app\par}
\vspace{2cm}

{\large\bfseries Final Report\par}
\vspace{1.5cm}

\begin{tabular}{rl}
\textbf{Author:}     & Angaraju V \\
\textbf{Team:}       & SLAP, Flipkart \\
\textbf{Mentor:}     & \textit{(mentor name)} \\
\textbf{Period:}     & May 2026 -- July 2026 \\
\textbf{Repository:} & \url{https://github.com/angarajuNiveditha/slap-bug-triage} \\
\textbf{Submitted:}  & July 2026 \\
\end{tabular}

\vfill

\begin{minipage}{0.85\textwidth}
\small\itshape
This report describes an end-to-end prototype I built to automate the first pass
of bug triage for Flipkart's SLAP (Shop Like A Pro) shopping app. The system
takes a raw bug report --- either as an email or through a structured form ---
optionally with screenshots or video, and produces a Jira-ready ticket draft
routed to the correct team, priority-scored, deduplicated against real
historical bugs, and assigned to a plausible owner. The system is read-only
against Jira by design (never writes) and always leaves the final call to a
human reviewer.
\end{minipage}

\end{titlepage}

% ────────────────────────── TOC ──────────────────────────
\tableofcontents
\newpage

% ────────────────────────── SECTION: EXECUTIVE SUMMARY ──────────────────────────
\section{Executive Summary}

SLAP is Flipkart's conversational shopping app --- a chat-first surface where a
user asks the assistant to find products, gets AI-generated recommendations,
tries clothes on virtually (VTON), browses generated Style Drops, and checks
out. Because the surface is heterogeneous (React Native shell, native
Android/iOS modules, backend chat AI over Edison, backend ML for VTON /
Styledrops / Social Finds, data-science layers for ranking / grounding), an
incoming bug report can plausibly belong to any of five teams. Triage is
currently manual and slow.

This project builds a prototype triage assistant that, given a bug report,
does the following:

\begin{enumerate}[nosep]
\item Parses text and (optionally) media attachments into a structured
      \texttt{BugReport}.
\item Classifies the owning team using an embedding-based Logistic Regression
      classifier, with a Claude fallback for borderline cases.
\item Retrieves the top-10 most similar historical bugs using a two-stage
      cosine-plus-cross-encoder retrieval.
\item Runs deduplication, owner suggestion and priority scoring in parallel.
\item Produces a Jira ticket draft (ADF JSON) and surfaces it in a Streamlit
      UI where a human reviewer can override any field before publishing.
\end{enumerate}

The system is intentionally \textbf{read-only against Jira} --- it never
files a ticket automatically. A human always confirms.

Three pipelines exist in the repository. The \textbf{multi-agent pipeline} is
the primary one and combines local ML with sub-agents backed by Anthropic's
Claude (via the Claude Code CLI, no API key needed) and Google's Gemini (via
Flipkart's internal proxy). The \textbf{rule-based pipeline} is a
fully-deterministic baseline (regex + TF-IDF) that runs in about 35 ms per
bug. The \textbf{Anthropic SDK pipeline} is the original design, kept as a
production-shape reference; it is currently blocked because
\texttt{console.anthropic.com} is unreachable on Flipkart's corporate network.

Measured accuracy of the component classifier: \textbf{69.5\% leave-one-out
accuracy} on 564 real labelled FLIPPI bugs. A label-noise audit found that
roughly 70\% of the classifier's mistakes are actually mis-labelled bugs in
Jira; after relabelling, projected accuracy is 78--82\% with no model
changes.

The rest of this document describes the design decisions, the code layout,
the accuracy work, and the exact steps to reproduce everything on a fresh
laptop.

\vspace{1em}
\begin{screenshotbox}[Landing screenshot of the Streamlit app]
Save the full-page screenshot of the app (with a bug in the input, no results
yet) at \texttt{report/screenshots/01\_landing.png} and it will render here.
Uncomment the \texttt{includegraphics} line below.

\begin{verbatim}
% \includegraphics[width=\linewidth]{screenshots/01_landing.png}
\end{verbatim}
\end{screenshotbox}

\newpage

% ────────────────────────── SECTION: PROBLEM ──────────────────────────
\section{The problem I was solving}

Bug triage on SLAP has three practical pain points:

\subsection{The team boundary is unclear}

A bug titled \textit{``the chat is broken''} could belong to any of:

\begin{itemize}[leftmargin=1.2em,nosep]
\item \textbf{Backend}: the chat API returned a 500 or timed out
\item \textbf{DS}: the model returned an irrelevant response
\item \textbf{UI}: the React Native chat screen crashed or froze
\item \textbf{Backend-Labs}: the chat was in a Styledrops / Social Finds context
\item \textbf{Immersive}: the failure was in a 3D / AR-VR view embedded in chat
\end{itemize}

Manual triage requires either domain knowledge or long back-and-forth on
Slack. The classifier I built handles the routine cases fast and flags the
genuinely ambiguous ones to a human.

\subsection{Historical bugs are a wasted signal}

Every bug that came before is a labelled example (component + assignee +
priority). Without automation, that signal is used at most as a mental
reference by a triage engineer who remembers the last similar bug. I wanted
to close that loop --- given a new bug, actually surface the top few similar
past bugs with their assignees and priorities, and use those to inform the
routing decision.

\subsection{Media adds another layer}

Bug reports often include a screenshot or a screen-recording. The reporter
might describe the bug incorrectly (e.g.\ ``checkout is broken'' when the
screenshot shows the cart page). Without automation, media has to be
manually inspected. I wanted the pipeline to look at the media, extract the
visible evidence, and use that to correct the reporter's framing when
necessary.

\subsection{What I was NOT trying to solve}

To keep the scope realistic:

\begin{itemize}[leftmargin=1.2em,nosep]
\item \textbf{Not auto-filing.} The system never writes to Jira. It produces
      a draft; a human clicks ``file'' after reviewing.
\item \textbf{Not full production.} This is a prototype. Real deployment
      would need write access to Jira via MCP, service-account authentication,
      metrics/monitoring and a scale-out story --- none of which is in scope
      here.
\item \textbf{Not model training from scratch.} I use pre-trained
      sentence-transformers and pre-trained cross-encoders. Only the
      Logistic Regression classifier is trained on our data, and that
      training is one-shot on 564 examples --- less than a second.
\end{itemize}

\newpage

% ────────────────────────── SECTION: ARCHITECTURE ──────────────────────────
\section{System architecture}

The system is organised into three layers: the \textbf{data layer} (real
FLIPPI bugs from Jira, indexed once and cached to disk), the
\textbf{pipeline layer} (three pipelines that share the data layer), and the
\textbf{presentation layer} (a Streamlit front-end). All three layers live
in a single Python codebase.

\begin{figure}[h]
\centering
\begin{minipage}{0.95\textwidth}
\begin{lstlisting}
                     ┌───────────────── DATA LAYER ─────────────────┐
                     │                                              │
                     │  Jira REST API (read-only)                   │
                     │      │                                       │
                     │      ▼                                       │
                     │  jira_client.py  ── 564 labelled bugs        │
                     │      │                                       │
                     │      ▼                                       │
                     │  build_embedding_index.py                    │
                     │      │                                       │
                     │      ▼                                       │
                     │  data/embedding_index.npz  (564x768 floats)  │
                     │  data/embedding_index_logreg.pkl             │
                     │  data/embedding_index_team_roster.json       │
                     │                                              │
                     └──────────────────────┬───────────────────────┘
                                            │
                                            │  loaded once at startup
                                            ▼
   ┌──────────────── PIPELINE LAYER (three, sharing the index) ────────────┐
   │                                                                       │
   │   [1] Rule-based   run_agent.py            ~35 ms per bug, no LLM     │
   │   [2] Multi-agent  run_multi_agent.py      ~40--70 s per bug, primary │
   │   [3] SDK          main.py                 blocked (no API key)       │
   │                                                                       │
   └───────────────────────────────┬───────────────────────────────────────┘
                                   │
                                   ▼
   ┌───────────────── PRESENTATION LAYER ────────────────────────┐
   │                                                             │
   │  Streamlit app (app.py)                                     │
   │  - input format toggle (email vs structured form)           │
   │  - pipeline toggle (rule-based vs multi-agent)              │
   │  - file uploader for screenshots / videos                   │
   │  - editable Priority / Component / Owner / Duplicate tiles  │
   │  - override -> data/corrections.csv (active learning)       │
   │                                                             │
   └─────────────────────────────────────────────────────────────┘
\end{lstlisting}
\end{minipage}
\caption{Three-layer architecture.}
\end{figure}

The data layer is populated once (whenever you rebuild the index --- takes
about a minute) and read-only after that. Both the rule-based and the
multi-agent pipelines share the same 564-bug corpus; the difference is what
they do with it.

\subsection{Three pipelines, three purposes}

\begin{table}[h]
\centering
\small
\begin{tabularx}{\linewidth}{lXll}
\toprule
\textbf{Pipeline} & \textbf{Purpose} & \textbf{Latency} & \textbf{Status} \\
\midrule
Rule-based           & Deterministic baseline. Regex + TF-IDF cosine over 300 bugs. No LLM anywhere. Used as the fallback when Claude / Gemini can't be reached, and as a control condition for measuring the multi-agent lift. & \textasciitilde 35 ms/bug & Working \\
\addlinespace
Multi-agent          & Primary path. Combines local ML (embedding-based classifier and similarity engine) with sub-agents backed by Claude (via the Claude Code CLI --- no API key required) and Gemini (via Flipkart's internal proxy). This is what the mentor will demo. & \textasciitilde 40--70 s/bug & Working \\
\addlinespace
Anthropic SDK        & Original design using the Anthropic Python SDK directly. Kept in the repo as a ``production-shape reference'' but blocked because \texttt{console.anthropic.com} is unreachable on Flipkart's corporate network. & --- & Blocked \\
\bottomrule
\end{tabularx}
\caption{The three pipelines in the repository.}
\end{table}

Everything below describes the \textbf{multi-agent pipeline} unless
explicitly noted, since that is the one the mentor will demo.

\newpage

% ────────────────────────── SECTION: DATA LAYER ──────────────────────────
\section{Data layer}

The system depends on a single cached artifact: an embedding index of 564
real labelled FLIPPI bugs pulled from Jira. Everything downstream reads from
this index.

\subsection{Fetching from Jira}

\texttt{src/jira\_client.py} wraps the Jira Cloud REST v3 API. Authentication
is HTTP Basic with an email + Atlassian PAT. The client has \textbf{no write
methods by design} --- it can search, get, and list, but not create, edit,
or transition. This is enforced at the code level, not just by convention.

\begin{lstlisting}[language=python,caption={jira\_client.py sample --- read-only surface}]
class JiraClient:
    def search_jql(self, jql: str, ...) -> list[dict]: ...
    def fetch_recent_bugs(self, limit: int) -> list[dict]: ...
    def get_issue(self, key: str) -> dict: ...
    # (no create_issue / edit_issue / transition_issue by design)
\end{lstlisting}

\subsection{Building the embedding index}

Once the raw Jira data is fetched, \texttt{build\_embedding\_index.py}
does the following, one-shot:

\begin{enumerate}[leftmargin=1.4em,nosep]
\item Filter to bugs that have a component label matching one of the five
      real teams (\texttt{Backend}, \texttt{Backend-Labs}, \texttt{DS},
      \texttt{UI}, \texttt{immersive}).
\item Concatenate \texttt{title + description} for each bug.
\item Encode each via \texttt{sentence-transformers/all-mpnet-base-v2}
      (768-dim, L2-normalised).
\item Train a Logistic Regression classifier on the resulting
      $(564, 768)$ feature matrix, with the team labels as ground truth.
\item Derive a team-roster mapping (component $\rightarrow$ engineers with
      $\geq 2$ bug assignments, managers excluded).
\item Cache all of the above to three files:
\end{enumerate}

\begin{lstlisting}
data/embedding_index.npz            (564x768 float32 + labels + keys)
data/embedding_index_logreg.pkl     (trained sklearn LogReg pickle)
data/embedding_index_team_roster.json  (component -> list of engineers)
\end{lstlisting}

These three files together are the shared state that all downstream
components read from. They are gitignored --- each contributor generates
their own copy from live Jira data.

\subsection{Why 564 bugs}

The number is not a design decision --- it's what the Jira filter returned
for ``recent labelled bugs from the FLIPPI project.'' The distribution is
uneven:

\begin{table}[h]
\centering
\small
\begin{tabular}{lrl}
\toprule
\textbf{Component} & \textbf{Bug count} & \textbf{Comment} \\
\midrule
Backend        & \textasciitilde 200 & Largest class \\
UI             & \textasciitilde 235 & Includes native + RN bugs \\
Backend-Labs   & \textasciitilde 120 & VTON, Styledrops, Social Finds, etc. \\
DS             & \textasciitilde 80  & Ranking, prompt, content quality \\
immersive      & 0                   & Zero bugs labelled ``immersive'' in the corpus \\
\bottomrule
\end{tabular}
\caption{Class distribution in the 564-bug training corpus.}
\end{table}

The uneven distribution is dealt with in the classifier via
\texttt{class\_weight="balanced"} (see \S\ref{sec:classifier}). The
zero-bugs-for-immersive problem is real and unresolved --- the classifier
has literally no idea what an immersive bug looks like.

\newpage

% ────────────────────────── SECTION: PIPELINE ──────────────────────────
\section{The multi-agent pipeline}\label{sec:pipeline}

\texttt{src/agents/host\_agent.py} is a thin coordinator (I called it
Astral, after the production agent-runtime term). It calls each stage in
sequence, threading a callback for progress updates so the Streamlit UI can
render a live status stepper. Independent stages run in parallel via
\texttt{concurrent.futures.ThreadPoolExecutor}.

\begin{figure}[h]
\centering
\begin{minipage}{0.95\textwidth}
\begin{lstlisting}
bug input (email OR structured form + optional images / videos)
    |
    v
[1] subagent_media           - only if attachments present     [Gemini + Claude]
    |     Stage A: Gemini describes each image / video keyframe
    |              in parallel via ThreadPoolExecutor
    |     Stage B: Claude reads the descriptions + SLAP knowledge
    |              + reference screens, emits MediaFinding JSON
    |     Falls back to legacy single-Claude vision if Gemini
    |     is unreachable (JWT expired, off-corp).
    v
[2] subagent_parser          - email + media -> BugReport     [Claude]
    |
    v
[2b] subagent_form_consistency - only if from_form=True       [Claude]
    |                            refile if title/summary/steps mismatch
    v
[3] EmbeddingClassifier      - LogReg (~7 ms).                [local ML]
    |    Confidence >= 0.60 -> return; else Claude+skills
    |    fallback with top-3 candidate teams' skill files      [Claude+skills]
    |
    v
[4] EmbeddingSimilarityEngine - TWO-STAGE:                    [local ML + rerank]
    |    (i)  cosine top-30 recall (~200 us)
    |    (ii) cross-encoder ms-marco-MiniLM-L-6-v2 rerank
    |         of the 30 pairs (~1.2 s on CPU) -> top-10
    v
[5] PARALLEL BLOCK (ThreadPoolExecutor, ~8-12 s wall clock)
    +- subagent_dedup        - dup/no-dup over top-10         [Claude]
    |                          (fires only if conf >= 0.80)
    +- subagent_owner        - engineer pick, else            [Claude + rules]
    |                          closest-similar-bug manager,
    |                          else TEAM_MANAGERS[component]
    +- subagent_triage       - priority P0 / P1 / P2          [Claude]
    v
agent_ticket_builder         - assembles Jira ADF + triage_notes JSON
    |
    v
human override (UI)          - reviewer edits Priority / Component / Owner;
                              audit trail in triage_notes.human_overrides
\end{lstlisting}
\end{minipage}
\caption{The full multi-agent pipeline. Times are on a mid-range CPU.}
\end{figure}

\subsection{Per-stage responsibilities}

\begin{table}[h]
\centering
\footnotesize
\begin{tabularx}{\linewidth}{lXl}
\toprule
\textbf{Stage} & \textbf{What it does} & \textbf{Implementation} \\
\midrule
Media & Reads attachments, identifies the SLAP screen, extracts visible bug evidence & Gemini vision + Claude reasoning \\
Parser & Turns the raw text into a structured \texttt{BugReport} & Claude \\
Form consistency & Only if from a structured form: flags title/summary/steps mismatches & Claude \\
Component classifier & Assigns the bug to one of five teams & LogReg + Claude fallback \\
Similarity engine & Ranks top-10 most similar past bugs & Cosine + cross-encoder rerank \\
Dedup & Decides whether the new bug is a duplicate of one of the top-10 & Claude (fires if conf $\geq 0.80$) \\
Owner & Picks the most likely engineer, escalates to team manager if needed & Claude + deterministic rules \\
Triage & Assigns P0 / P1 / P2 with a plain-English justification & Claude \\
Ticket builder & Assembles the final ADF JSON draft & Python \\
\bottomrule
\end{tabularx}
\caption{Per-stage responsibilities in the multi-agent pipeline.}
\end{table}

\begin{screenshotbox}[Streamlit live-progress stepper during a triage run]
Save a screenshot of the pipeline actually running (with the nested
progress events visible: \texttt{media:vision:start},
\texttt{media:reasoning:done}, etc.) at
\texttt{report/screenshots/02\_pipeline\_running.png}.
\end{screenshotbox}

\newpage

% ────────────────────────── SECTION: MEDIA ──────────────────────────
\section{Media sub-agent (Gemini vision + Claude reasoning)}

Media is where the biggest engineering decision landed. The original design
had Claude read image attachments directly (via the \texttt{--add-dir} flag
that grants the CLI Read-tool access to the attachment folder). That worked
but was slow: about 12--15 seconds per image, and serial across multiple
images.

I split it into two stages:

\begin{itemize}[leftmargin=1.4em,nosep]
\item \textbf{Stage A --- Gemini (vision)}: Each image (or each of the
      $\leq 8$ ffmpeg-extracted video keyframes) is described by Gemini
      2.5 Flash in \emph{parallel} via a \texttt{ThreadPoolExecutor}.
      Gemini just describes what's visible; it doesn't know anything about
      SLAP semantics. About 3 s wall clock for images regardless of count;
      about 14 s for 8 video keyframes.
\item \textbf{Stage B --- Claude (reasoning)}: A single Claude call reads
      the plain-prose Gemini descriptions plus
      \texttt{slap\_context/SLAP\_KNOWLEDGE.md} (a Figma-extracted screen
      catalog) and \texttt{slap\_context/reference\_screens/} (16 labelled
      Figma PNGs), and emits the structured \texttt{MediaFinding} JSON.
      Claude does the SLAP-specific reasoning (screen mapping, triage
      signals, contradiction detection between email claim and visual
      evidence).
\end{itemize}

If Gemini is unreachable (JWT expired, off-corp) the sub-agent falls back
to the legacy single-Claude-call path. The pipeline never hard-fails on
media problems.

\begin{decisionbox}[Use Gemini for vision, Claude for SLAP reasoning]
Gemini is faster at raw vision (parallel across images / frames) and
requires no SLAP context. Claude is better at multi-step reasoning over
structured input and already has SLAP context loaded. Splitting the two
lets each model do what it's best at. Fallback path preserved for
network / auth failures.
\end{decisionbox}

\subsection{The Gemini proxy plumbing}

Gemini isn't reached directly. Flipkart routes Gemini through an internal
proxy at \texttt{http://10.83.64.112/gemini-2.5-flash/:generateContent}.
The endpoint requires \emph{both} an Azure APIM subscription key
(\texttt{Ocp-Apim-Subscription-Key} header) \emph{and} a short-lived JWT
bearer (\texttt{Authorization: Bearer <GENVOY\_TOKEN>}). Neither alone
works.

The dual-auth combination isn't documented anywhere I could find --- I
figured it out by probing (sending different combinations of headers to
the endpoint and reading the error messages, which spelled out what was
missing).

\begin{gotchabox}[The Genvoy JWT expires every hour]
The bearer token has a one-hour TTL. When it lapses mid-session, the
Streamlit app used to keep running against an expired token because
Python-dotenv doesn't override existing env vars on reload. I added
\texttt{load\_dotenv(override=True)} so that editing \texttt{.env} and
hitting any Streamlit widget picks up the new value, plus a visible
banner in the UI that goes red when the token is expired.
\end{gotchabox}

\subsection{Videos: keyframe extraction + sequence reasoning}

Videos are pre-processed with \texttt{imageio-ffmpeg} (portable ffmpeg
binary, no system ffmpeg required). Scene-detection extracts up to 8
keyframes; if the scene detector returns nothing (very static video),
evenly-spaced sampling is used as a fallback. Videos longer than 60 s are
rejected with a clear message asking the reporter to trim.

Once the keyframes exist, the same two-stage flow runs: Gemini describes
each frame in parallel, Claude reasons over the sequence and produces one
\texttt{MediaFinding} per video (with \texttt{screen\_sequence},
\texttt{action\_observed}, \texttt{failure\_moment} fields populated).

\begin{screenshotbox}[Media findings JSON for a real bug]
Save a screenshot of the Streamlit ``Findings'' tab showing the full
Media finding structure at
\texttt{report/screenshots/03\_media\_findings.png}. The
\texttt{contradicts\_email\_claim} field is worth highlighting --- it's
often the most useful field for a triage engineer.
\end{screenshotbox}

\newpage

% ────────────────────────── SECTION: CLASSIFIER ──────────────────────────
\section{Component classifier}\label{sec:classifier}

The classifier answers ``which team owns this?'' It runs on the same 768-dim
sentence-transformer embedding used elsewhere.

\subsection{Why Logistic Regression}

I considered three options:

\begin{table}[h]
\centering
\small
\begin{tabularx}{\linewidth}{lXl}
\toprule
\textbf{Approach} & \textbf{Reason not / reason yes} & \textbf{Verdict} \\
\midrule
k-Nearest Neighbours & No probability output, so can't drive a confidence-based fallback. Susceptible to label noise --- one mis-labelled Backend bug in the top-5 flips the vote. & Rejected \\
\addlinespace
Small MLP           & 564 examples are far too few to train even a small MLP without overfitting. And a linear model on top of a good embedding is usually the right tool at this data scale. & Rejected \\
\addlinespace
Logistic Regression & Fast (7 ms per prediction), outputs a probability distribution that we can threshold, hard to overfit at $5\times 768=3{,}840$ parameters, well-suited to already-embedded feature spaces. & \textbf{Chosen} \\
\bottomrule
\end{tabularx}
\end{table}

\subsection{Training and class imbalance}

Training is one-shot inside \texttt{build\_embedding\_index.py}:

\begin{lstlisting}[language=python]
logreg = LogisticRegression(
    class_weight = "balanced",  # counter the ~5:1 imbalance
    max_iter     = 2000,
    C            = 1.0,         # standard L2 regularisation
)
logreg.fit(embeddings, labels)  # (564, 768) @ 5 classes
\end{lstlisting}

\texttt{class\_weight="balanced"} counters the fact that Backend and UI
together are $\sim$65\% of the corpus. Without it, LogReg would learn
``predict Backend when uncertain'' and the smaller classes would be
under-represented in the decision boundary.

\subsection{The hybrid fast path}

At prediction time the classifier applies a hybrid rule:

\begin{lstlisting}[language=python]
HYBRID_CLAUDE_FALLBACK_THRESHOLD = 0.60

def predict(self, text: str) -> ClassificationResult:
    q = self.embedder.encode([text], normalize_embeddings=True)[0]
    proba = self.logreg.predict_proba(q.reshape(1, -1))[0]
    if max(proba) >= HYBRID_CLAUDE_FALLBACK_THRESHOLD:
        return LogReg's argmax        # fast path (~7 ms)
    else:
        return Claude+skills fallback  # slow path (~6 s), passes top-3
                                       # candidate teams' skill files
\end{lstlisting}

The threshold was originally set at 0.50 and later raised to 0.60. The
higher threshold pushes more borderline cases to Claude+skills, trading
about 600 ms of average latency for a small accuracy gain on the shifted
bugs.

\begin{decisionbox}[Threshold 0.60 for the fast/slow-path switch]
Bugs where LogReg's top-class probability is between 0.50 and 0.60 are
``somewhat confident'' picks that historically have lower accuracy than
$\geq 0.60$ picks. Routing them to Claude+skills lets Claude reason
over architecture context (skill files) rather than trusting LogReg's
partial confidence. The cost is $\sim$600 ms per bug on average.
\end{decisionbox}

\subsection{Measured accuracy}

Leave-one-out cross-validation on the 564-bug corpus:

\begin{table}[h]
\centering
\small
\begin{tabular}{lr}
\toprule
\textbf{Approach} & \textbf{LOO Accuracy} \\
\midrule
Old keyword-regex classifier                & 27.7\% \\
Pure LogReg (no fallback)                   & 66.8\% \\
Hybrid LogReg + Claude+skills @ threshold 0.50 & \textbf{69.5\%} \\
Projected after backend label cleanup       & 78--82\% (est.) \\
\bottomrule
\end{tabular}
\caption{Classifier accuracy across iterations.}
\end{table}

The threshold-0.60 variant has not yet been re-benchmarked through the LOO
harness --- I expect a small uptick on the shifted bugs but haven't
measured. This is a low-effort follow-up.

\subsection{The label-noise finding}

An audit of the misclassifications showed something interesting: of the
$\sim$60 misclassified Backend bugs, roughly 70\% were \emph{mis-labelled}
in Jira. Chat-AI / model-quality complaints (``bot didn't understand'',
``wrong results'', ``reasons to buy phrased oddly'') were being filed
against the \texttt{Backend} component but are textbook DS bugs. The
classifier correctly identifies them as DS-shaped; we're scoring it against
noisy ground truth.

If those 60 bugs were relabelled, projected accuracy jumps to 78--82\% with
no model changes. This is the single highest-leverage improvement
available and is documented as ``future work.''

\newpage

% ────────────────────────── SECTION: RETRIEVAL ──────────────────────────
\section{Similar-bug retrieval (cosine + cross-encoder)}

Retrieval answers ``which past bugs are like this one?'' It runs in two
stages.

\subsection{Stage 1: cosine top-30 recall}

Both the query text and each of the 564 past bugs are represented as
L2-normalised 768-dim vectors. Similarity is just the dot product:

\begin{lstlisting}[language=python]
sims    = self.embeddings @ q   # (564, 768) @ (768,) -> (564,)
top_idx = np.argpartition(-sims, 30)[:30]
\end{lstlisting}

About 200 microseconds for 564 candidates. High recall but fuzzy at fine
ranking --- cosine cares only about geometric proximity of the two
embeddings, which were produced independently and don't ``see'' each
other.

\subsection{Stage 2: cross-encoder rerank of the top-30}

The 30 candidates are then rescored by
\texttt{cross-encoder/ms-marco-MiniLM-L-6-v2}, which is a small BERT-family
model (22M parameters) that takes \emph{two} texts as input and produces a
single relevance score. Because both texts go through the model together,
the attention layers can compare tokens across the two texts directly
--- something a bi-encoder can't do.

\begin{decisionbox}[Cross-encoder rerank of top-30, not raw cross-encoder retrieval]
Running the cross-encoder against all 564 candidates would take 5--10 s
per query. Running it against only the top-30 cosine candidates costs
about 1.2 s. The extra quality of ``cross-encode everything'' over
``cross-encode the cosine shortlist'' is essentially zero, because the
true top-K is almost always inside the cosine top-30.
\end{decisionbox}

\subsection{Example: why this matters}

Test bug: \texttt{[Checkout]: App crashes on Proceed to Pay, Android 2.4.2}.

\begin{table}[h]
\centering
\footnotesize
\begin{tabular}{clcc}
\toprule
\textbf{Rank} & \textbf{Bug} & \textbf{Cosine only} & \textbf{After rerank} \\
\midrule
1 & FLIPPI-616 (``Buy Now'' tap crash) & 0.685 & 0.411 \\
2 & FLIPPI-1198 (iOS ``continue to payment'' crash) & 0.624 & \textbf{0.790} \\
3 & FLIPPI-774 (crash while switching apps) & 0.582 & 0.508 \\
\bottomrule
\end{tabular}
\caption{Ranking flip: cosine's top pick (FLIPPI-616, matched on ``Buy Now button'' surface tokens) gets demoted, and the cross-encoder correctly promotes FLIPPI-1198 (iOS parallel of the same failure mode) to \#1.}
\end{table}

The cross-encoder ``understood'' that ``Proceed to Pay'' and ``continue to
payment'' are the same failure moment even though the words are different.
The bi-encoder could not, because it produced the two embeddings
independently and lost that alignment.

\subsection{Score normalisation}

Raw cross-encoder scores don't have a fixed range (unlike cosine's
$[-1, 1]$). I sigmoid-normalise them to $[0, 1]$ before returning, so
downstream code (dedup / owner) can compare them monotonically. The
\texttt{is\_duplicate\_candidate} flag on each \texttt{SimilarBug} still
tracks the raw cosine value, because the 0.80 duplicate threshold is
calibrated for the cosine scale.

\subsection{Relevance gate: don't reason over noise}\label{sec:relevance-gate}

Sometimes the corpus genuinely has nothing similar to a new bug. When
that happens the cross-encoder produces a top-10 with sigmoid scores in
the 0.01--0.03 band --- i.e.\ raw scores around $-4$, which is the
model saying ``these are confidently not related.'' If we still hand
that list to owner and triage, they end up reasoning over unrelated
bugs: owner picks from assignees who worked on irrelevant tickets,
triage averages priorities of irrelevant tickets.

To stop that, the retrieval carries a low-confidence flag:

\begin{lstlisting}[language=python]
RELEVANCE_THRESHOLD = 0.10   # sigmoid-normalised CE score

@dataclass
class EmbeddingSimilarityResult:
    top_matches:         list       # always returned to the UI
    top_relevance_score: float      # max .similarity across top_matches
    is_low_confidence:   bool       # top_relevance_score < RELEVANCE_THRESHOLD
\end{lstlisting}

When \texttt{is\_low\_confidence} is true, the host passes an
\emph{empty} list to the parallel block instead of \texttt{top\_matches}.
Each downstream sub-agent already handles the empty case gracefully:

\begin{itemize}[leftmargin=1.4em,nosep]
\item \textbf{Dedup}: returns ``no duplicate'' trivially.
\item \textbf{Owner}: \texttt{same\_component} is empty $\rightarrow$
      cascades to \texttt{\_escalate\_to\_team\_manager(component)}.
\item \textbf{Triage}: sees an empty ``similar bugs'' block in the prompt
      and reasons from the bug text alone (the priority ladder + hard
      overrides in \texttt{subagent\_triage.py} still apply).
\end{itemize}

The full top-10 is still stored on the draft so the UI can show it to
the reviewer for context, with a warning banner above the table
explaining \emph{why} the downstream picks look like escalations.

\begin{decisionbox}[Threshold value 0.10, chosen for the calibration property, not empirically tuned]
$\text{sigmoid}(-2.2) \approx 0.10$. Below that, the raw cross-encoder
score is meaningfully negative --- the model is leaning ``not
related.'' Good matches typically score $0.6+$; loosely-related bugs
in the $0.20$--$0.40$ band; the $<0.10$ zone is where the model is
confidently ``not similar.'' If the demo turns up cases where this
feels too conservative, the constant is one line to change.
\end{decisionbox}

\newpage

% ────────────────────────── SECTION: DEDUP / OWNER / TRIAGE ──────────────────────────
\section{Dedup, owner, and triage --- in parallel}

Three sub-agents run in parallel after the classifier and similarity stages
finish:

\begin{lstlisting}[language=python]
with ThreadPoolExecutor(max_workers=3) as ex:
    fut_dedup   = ex.submit(decide_duplicate,   bug, top_matches)
    fut_owner   = ex.submit(suggest_owner,       bug, top_matches, roster)
    fut_triage  = ex.submit(score_severity,      bug, top_matches)
    dup, owner, severity = ...
\end{lstlisting}

Wall clock for this block drops from about 20--30 s (serial) to 8--12 s
(parallel), which is where most of the pipeline speedup came from.

\subsection{Dedup}

\texttt{subagent\_dedup.py} takes the top-10 similar bugs and asks Claude:
``Is the new bug a duplicate of one of these?'' Claude is told to fire only
when it's $\geq 0.80$ confident. False positives are worse than false
negatives here --- linking two bugs incorrectly wastes engineering time.

\subsection{Owner --- the escalation ladder}

This is where a meaningful design decision landed. The naive rule is ``pick
the most frequent assignee for this component.'' That rule fails because:

\begin{itemize}[leftmargin=1.4em,nosep]
\item Managers accidentally get counted (Yatin Grover appears as historical
      assignee on many bugs because he inherited them during triage, but he
      shouldn't be the default owner).
\item Engineers with lots of bugs on the \emph{component} but zero bugs on
      the \emph{specific failure area} would still win, which produces
      obviously-wrong picks.
\item Engineers whose primary team is different but who touched a single
      same-component bug during a rotation would surface as candidates too.
      (Real case: Kalpana lnm has 30 UI-labelled bugs and 1 stray
      Backend-labelled assignment. For a Backend-routed bug she was being
      suggested over Divya Chhibber, an actual Backend engineer.)
\end{itemize}

The rule I settled on, in order:

\begin{enumerate}[leftmargin=1.4em,nosep]
\item If any \emph{engineer} (non-manager) has been assigned a similar bug
      on this component \emph{AND is in the component's current roster}
      $\rightarrow$ Claude picks between them. The roster-membership check
      is what stops the Kalpana-type case.
\item If only \emph{managers} appear in the similar-bug assignees $\rightarrow$
      pick the manager who was on the single \emph{closest} similar bug
      (highest cross-encoder score).
\item If no similar bugs exist on this component at all $\rightarrow$ escalate
      to the component's manager per \texttt{TEAM\_MANAGERS} in
      \texttt{src/team\_config.py}:
      \begin{itemize}[nosep]
      \item UI, Backend-Labs, Immersive $\rightarrow$ Yatin Grover
      \item Backend $\rightarrow$ Veeramreddy ChakradharReddy
      \item DS $\rightarrow$ (not mapped --- no manager in the org chart yet;
            returns ``no candidates'' for manual triage)
      \end{itemize}
\end{enumerate}

The team roster itself comes from \texttt{data/embedding\_index\_team\_roster.json},
which is derived at index-build time from historical Jira assignees, filtered
to people with at least 2 bugs per component, and with \texttt{MANAGER\_NAMES}
stripped out. So a candidate for a Backend-routed bug is (a) non-manager AND
(b) has $\geq 2$ Backend bugs in history AND (c) shows up in the similar-bug
list for this specific bug. All three filters have to pass.

\begin{decisionbox}[Manager mapping is a one-line-of-config source of truth]
\texttt{src/team\_config.py} defines both \texttt{MANAGER\_NAMES} (excluded
from the derived roster during index build) and \texttt{TEAM\_MANAGERS}
(the component $\rightarrow$ manager map). Both are imported by
\texttt{app.py} (for the Owner dropdown label) and
\texttt{embedding\_classifier.py} (roster derivation) and
\texttt{subagent\_owner.py} (escalation), so there's one source of truth
and no drift.
\end{decisionbox}

\begin{gotchabox}[Regression that took two commits to notice]
The roster-membership filter existed in the very first version of the
owner sub-agent (June commit \texttt{653f2cb}). A later refactor that
added the manager-escalation ladder (\texttt{251c878}) inadvertently
changed a set-intersection to a set-union and silently broke the
filter --- I noticed only when a screenshot showed a UI engineer being
suggested for a Backend bug. Re-added in \texttt{7450b24}. A proper
regression harness against the ``candidate must be in the routed
team's roster'' invariant would have caught this in CI.
\end{gotchabox}

\subsection{Triage}

\texttt{subagent\_triage.py} takes the \texttt{BugReport} and the top-10
similar bugs and returns a priority (P0 / P1 / P2) with a justification. It
uses similar-bug priorities as the primary signal (``if the last three
similar bugs were P1, this is probably a P1 too'') and applies hard
overrides for real crashes (100\% reproducible crash $\rightarrow$ P0) and
security signals (Grayskull, auth infra $\rightarrow$ P0).

The ladder is three-tier by design (P0 / P1 / P2). I removed the P3 tier
entirely; if the classifier can't decide anything, the bug is routed to a
manual-triage queue instead of being labelled P3.

\newpage

% ────────────────────────── SECTION: SKILL FILES ──────────────────────────
\section{Architecture skill files}

When the classifier's fast path is uncertain (LogReg confidence $<0.60$),
Claude gets called with skill files loaded into the prompt. These files
describe what each team owns, the disambiguation rules between adjacent
teams, and the module structure inside each team's repos.

There are two levels:

\begin{itemize}[leftmargin=1.4em,nosep]
\item \textbf{Team-level skill files} at \texttt{slap\_context/architecture/}:
      one per team ({\tt Backend.md}, {\tt UI.md}, {\tt Backend-Labs.md},
      {\tt DS.md}, {\tt immersive.md}).
\item \textbf{Per-repo skill files} at
      \texttt{slap\_context/architecture/repos/}: one per repo
      ({\tt edison.md}, {\tt dropsense.md}, {\tt FaceNet.md},
      {\tt slap-feed.md}, {\tt social-finds-pipeline.md},
      {\tt slap-auto-qc-pipeline.md}, plus hand-authored
      {\tt spaghetti.md} and {\tt mozzarella.md}).
\end{itemize}

The per-repo files are \emph{auto-generated} from the actual cloned repos
by \texttt{build\_repo\_skills.py}. The script grep-mines each repo for:

\begin{itemize}[leftmargin=1.4em,nosep]
\item Service class names (\texttt{*Service.java} filenames)
\item Exception class names (\texttt{*Exception.java})
\item HTTP entry-point classes (\texttt{*Controller.java},
      \texttt{*Handler.java}, \texttt{*Endpoint.java}, \texttt{*Resource.java})
\item Spring routes (\texttt{@GetMapping}, \texttt{@PostMapping}, etc.)
\item Enum names (grep \texttt{public enum})
\item DTO / Request / Response class counts
\item Configuration files (\texttt{application*.yml/properties})
\end{itemize}

The result is that every claim in the skill files is grep-able against the
real code. When Claude reads
``\texttt{authentication/} contains \texttt{AuthenticationService},
\texttt{KevlarTokenService}, \texttt{OtpService}, \texttt{TokenService},''
that isn't a Claude hallucination --- those are real filenames in the
Edison repo.

\begin{decisionbox}[Mine the code, don't paraphrase from CLAUDE.md]
The first version of the team-level skill files was Claude-drafted prose
from CLAUDE.md, and the classifier's Claude fallback reasoning was
plausible-sounding but sometimes wrong (invented class names, wrong
module boundaries). Regenerating from mined code names produced a
demonstrable improvement in Claude's reasoning quality on borderline
bugs.
\end{decisionbox}

\subsection{Skill file coverage today}

\begin{table}[h]
\centering
\small
\begin{tabular}{lrl}
\toprule
\textbf{File} & \textbf{Lines} & \textbf{Source} \\
\midrule
\multicolumn{3}{l}{\textit{Team-level (hand-written on top of mined per-repo files)}} \\
\quad Backend.md          & 214 & rewritten 2026-07-07 \\
\quad UI.md               & 148 & rewritten 2026-07-07 \\
\quad Backend-Labs.md     & 147 & rewritten 2026-07-07 \\
\quad DS.md               & 142 & rewritten 2026-07-07 \\
\quad immersive.md        & 141 & org-chart only (no repo cloned) \\
\midrule
\multicolumn{3}{l}{\textit{Per-repo (auto-generated from clones)}} \\
\quad edison.md           & 170 & mined \\
\quad slap-feed.md        & 177 & mined \\
\quad dropsense.md        & 152 & mined \\
\quad social-finds-pipeline.md & 134 & mined \\
\quad FaceNet.md          & 123 & mined \\
\quad slap-auto-qc-pipeline.md & 78 & mined \\
\midrule
\multicolumn{3}{l}{\textit{Per-repo (hand-authored --- kept separately)}} \\
\quad spaghetti.md        & 255 & hand-authored \\
\quad mozzarella.md       & 284 & hand-authored \\
\bottomrule
\end{tabular}
\caption{Skill-file corpus. Total $\sim$105 KB across 13 files.}
\end{table}

\texttt{HAND\_AUTHORED = \{"spaghetti", "mozzarella"\}} in
\texttt{build\_repo\_skills.py} makes sure bulk regenerations don't
overwrite the two hand-authored files. Passing the name explicitly on the
command line still regenerates.

\newpage

% ────────────────────────── SECTION: UI ──────────────────────────
\section{The Streamlit UI (human-in-the-loop)}

The front-end is a single-file Streamlit app (\texttt{app.py}). It intentionally
keeps the human in the loop --- the pipeline produces a draft, the reviewer
edits any field that looks wrong, and the ``Publish to Jira'' button is a
no-op in the prototype (real deployment would wire it to the write side of
the Jira MCP).

Key UI features:

\begin{itemize}[leftmargin=1.4em,nosep]
\item \textbf{Input format toggle}: paste an email, or fill a structured
      form (title / platform / summary / steps). The form path also
      triggers the form-consistency sub-agent.
\item \textbf{Pipeline toggle}: rule-based vs multi-agent. The rule-based
      path is instant and text-only; the multi-agent path is slower but
      supports media.
\item \textbf{File uploader}: PNG / JPG / MP4 / MOV. Multi-agent only.
\item \textbf{Live progress stepper}: the pipeline emits events
      (\texttt{media:vision:start}, \texttt{media:reasoning:done},
      \texttt{similarity:done}, \texttt{parallel:start},
      \texttt{dedup:done}, \texttt{owner:done}, \texttt{triage:done}\ldots).
      Each event renders as a bullet with an icon and nested indentation
      matching the pipeline structure.
\item \textbf{Editable Step 2 tiles}: Priority / Component / Owner /
      Duplicate-of. Each is a selectbox (except Duplicate which is
      read-only) with the model's prediction as the default. Overrides go
      into \texttt{data/corrections.csv} for active learning on the next
      index rebuild.
\item \textbf{Ambiguity banner}: shown when LogReg confidence $<0.60$.
      Displays the full 5-class probability distribution.
\item \textbf{Token-expiry banner}: red at the top of the page when
      GENVOY\_TOKEN is expired / missing / malformed. Silent when
      healthy. Backed by \texttt{genvoy\_token\_status()} which decodes
      the JWT's \texttt{exp} claim and picks one of five states.
\item \textbf{Low-confidence retrieval warning}: a $\square$ banner
      above the Similar Past Bugs table when the top similarity was
      below the 0.10 relevance threshold. Tells the reviewer that
      owner + triage were escalated to defaults rather than derived
      from the shown bugs, so they don't misread the pipeline's
      escalation as ``no owner data available.''
\item \textbf{Widget-key versioning}: each fresh triage bumps a
      \texttt{triage\_version} counter and the Step 2 widget keys
      include it, so the tiles always reset to the newest predictions
      instead of persisting the previous override.
\item \textbf{Hot-reload of \texttt{.env}}: \texttt{load\_dotenv(override=True)}
      at module load, so refreshing the hourly-expiring \texttt{GENVOY\_TOKEN}
      in \texttt{.env} takes effect on the very next widget click. No
      Streamlit restart required.
\end{itemize}

\begin{screenshotbox}[Streamlit UI --- full page]
Save at \texttt{report/screenshots/04\_ui\_overview.png}. Ideally include
the pipeline rail on the left, the input on the right, and the editable
Step 2 tiles below.
\end{screenshotbox}

\begin{screenshotbox}[Owner dropdown with the manager escalation visible]
Save at \texttt{report/screenshots/05\_owner\_dropdown.png}. Show the
dropdown open with ``Yatin Grover $\cdot$ Manager'' and ``Veeramreddy
ChakradharReddy $\cdot$ Manager'' pinned at the top.
\end{screenshotbox}

\newpage

% ────────────────────────── SECTION: KEY DECISIONS ──────────────────────────
\section{Key design decisions --- a running record}

This section collects the decisions that mattered, with the reasoning I
recall for each.

\subsection{Read-only Jira}

The system never writes to Jira. This is enforced at the code level in
\texttt{jira\_client.py} (no create / edit / transition methods). The
``Publish to Jira'' button in the UI is a no-op. This choice keeps the
prototype defensible against a lot of concerns --- accidentally filing
duplicates, mis-routing, RBAC issues --- and defers the write-path work
to a real production integration.

\subsection{Three pipelines, one repo}

Both the rule-based baseline and the multi-agent pipeline live in the same
repo, sharing the data layer. The rule-based one runs in about 35 ms per
bug and is what the UI uses when the reviewer wants a quick sanity check
or when Claude / Gemini are unreachable. The multi-agent one is the demo
target. Keeping both meant I could always A/B them against each other,
which was invaluable when tuning the similarity engine.

\subsection{Claude Code CLI instead of the Anthropic SDK}

The Anthropic Python SDK needs an API key from
\texttt{console.anthropic.com}, which is unreachable on Flipkart's
corporate network. Instead, every ``ask Claude'' call in the multi-agent
pipeline goes through the Claude Code CLI (\texttt{claude -p}) via a small
subprocess wrapper (\texttt{src/claude\_cli.py}). Authentication is
inherited from the local Claude Code session --- no key needed. This is
the single choice that made the multi-agent pipeline runnable inside
Flipkart at all.

\subsection{Gemini via the internal proxy, not Google AI directly}

Google's public AI endpoint (\texttt{generativelanguage.googleapis.com})
is likewise unreachable. Flipkart has an internal proxy that fronts Gemini
with Azure APIM auth. I built a thin client (\texttt{src/genvoy\_client.py})
that handles the dual-header authentication and surfaces
\texttt{GeminiUnavailable} so callers can fall back gracefully.

\subsection{Two-stage retrieval instead of one-stage cross-encoder}

A pure cross-encoder retrieval would produce marginally better rankings
but at 5--10 s per query --- unusable. Cosine top-30 as the recall stage
plus a cross-encoder rerank of only those 30 candidates gets essentially
the same quality at $\sim$1.2 s. This is the standard modern retrieval
pattern and works well here.

\subsection{Parallel dedup / owner / triage}

These three sub-agents are independent (none reads the output of another).
Running them in a \texttt{ThreadPoolExecutor} saved 15--20 s per bug --- a
noticeable improvement for what turned out to be a 3-line change.

\subsection{The escalation ladder for owner picking}

The obvious ``most-frequent assignee'' rule kept picking wrong people
(managers who inherited old bugs; engineers with lots of unrelated bugs
on the same component). The current rule
(``similar-bug engineer $>$ closest-similar-bug manager $>$
team manager'') gets the right answer on the specific cases the naive
rule was failing on, and it degrades gracefully to a team manager when
there's nothing else to go on.

\subsection{Skill files grounded in mined code}

The original team-level skill files were Claude-drafted from CLAUDE.md.
Regenerating with real class names / exception names / enum values / route
paths from the cloned repos produced a demonstrable improvement in
Claude's fallback classification quality. Grep-verifiable prose is worth
the code that generates it.

\subsection{The 0.60 confidence threshold}

Raised from an original 0.50 to 0.60 to route more borderline bugs to
Claude+skills. This costs about 600 ms of average latency in exchange for
a small accuracy bump on the shifted bugs. The trade-off is a defensible
one for a demo where accuracy matters more than latency.

\subsection{Enforcing team-roster membership on owner candidates}

The owner sub-agent's candidate pool is now the intersection of ``non-
manager assignees on same-component similar bugs'' AND ``members of the
current roster for the routed component.'' Both filters have to pass ---
being on a same-component similar bug alone isn't enough if the person's
primary team is different. The alternative (union-based candidate pool,
which the code briefly had) let a UI engineer be suggested for a Backend
bug because they had one stray Backend assignment in history. Making the
filter an intersection restored the property that classification-time
routing (``this is a Backend bug'') strictly implies pick-time routing
(``the owner must be a Backend engineer''). The two are otherwise
allowed to disagree, and every disagreement is a bug.

\subsection{The 0.10 relevance-gate threshold on retrieval}

Cross-encoder scores in the $0.01$--$0.03$ band are the model
confidently saying ``these bugs are not related.'' Handing that list
to owner and triage produces plausible-looking-but-wrong picks (owner
picks an engineer who worked on an unrelated bug; triage averages
priorities of unrelated bugs). The relevance gate short-circuits that:
if the top rerank score is below \texttt{RELEVANCE\_THRESHOLD = 0.10},
the host passes an empty list to the reasoning sub-agents, and each
sub-agent's ``no similar bugs'' path handles it correctly (owner
escalates, triage reasons from bug text, dedup returns no duplicate).
Details in section~\ref{sec:relevance-gate}.

\newpage

% ────────────────────────── SECTION: REPRODUCTION ──────────────────────────
\section{Reproduction steps --- from a clean laptop}

These steps assume a mid-2020s Mac (Intel or Apple Silicon), Python 3.11+,
and a working Homebrew install. Linux is very similar; Windows is
untested.

\subsection{Step 1: Clone the repository}

\begin{lstlisting}[language=bash]
git clone https://github.com/angarajuNiveditha/slap-bug-triage.git
cd slap-bug-triage
\end{lstlisting}

\subsection{Step 2: Install and sign in to Claude Code}

\begin{lstlisting}[language=bash]
brew install --cask claude
# Then open Claude Desktop and sign in with your Flipkart Google account
\end{lstlisting}

Verify:

\begin{lstlisting}[language=bash]
which claude          # should print a path
claude -p "hi"        # should respond with a short reply
\end{lstlisting}

\subsection{Step 3: Python dependencies}

\begin{lstlisting}[language=bash]
python3 -m venv myenv
source myenv/bin/activate
pip install -r requirements.txt
\end{lstlisting}

\subsection{Step 4: Configure \texttt{.env}}

\begin{lstlisting}[language=bash]
cp .env.example .env
# then edit .env
\end{lstlisting}

Values to set:

\begin{lstlisting}
JIRA_EMAIL=<your flipkart email>
JIRA_TOKEN=<your atlassian PAT>
JIRA_BASE_URL=https://flipkart.atlassian.net
JIRA_PROJECT=FLIPPI

ANTHROPIC_API_KEY=        # leave empty
CLAUDE_BIN=               # leave empty; auto-detected on most Macs

# Gemini via Flipkart's internal proxy (for the fast media path)
GEMINI_API_URL=http://10.83.64.112/gemini-2.5-flash/:generateContent
GEMINI_API_KEY=<APIM subscription key>
GENVOY_BASE_URL=http://genvoy.jarvis-prod.fkcloud.in
GENVOY_SUBSCRIPTION_KEY=<same as GEMINI_API_KEY>
GENVOY_TOKEN=<short-lived JWT bearer; expires hourly>
\end{lstlisting}

\begin{gotchabox}[GENVOY\_TOKEN expires hourly]
The bearer JWT has a one-hour TTL. When it lapses, the pipeline still
runs but the media sub-agent falls back to Claude-only vision (slower).
The Streamlit UI shows a red banner at the top of the page when the
token is expired --- refresh it in \texttt{.env} and any click will pick
up the new value.
\end{gotchabox}

\subsection{Step 5: Build the embedding index (one-time)}

\begin{lstlisting}[language=bash]
python3 build_embedding_index.py
\end{lstlisting}

This fetches $\sim$564 component-labelled FLIPPI bugs from Jira, downloads
the sentence-transformer model on first run ($\sim$420 MB, one-time),
trains the LogReg classifier, derives the team roster, and writes three
files to \texttt{data/}. Takes about one minute total after the model
download.

\subsection{Step 6: (Optional) Build per-repo skill files}

If you have Flipkart GitHub Enterprise access and want to regenerate the
per-repo skill files from live clones:

\begin{lstlisting}[language=bash]
# Clone the SLAP repos into data/repos/<name>/ manually, then:
python3 build_repo_skills.py
\end{lstlisting}

This step is optional. The repository already has all skill files
committed; you only need to regenerate if the code changed materially.

\subsection{Step 7: Run the Streamlit UI}

\begin{lstlisting}[language=bash]
streamlit run app.py
# Opens at http://localhost:8501
\end{lstlisting}

Try one of the sample bugs (dropdown in the input area picks up any
\texttt{.txt} file under \texttt{data/}). Toggle between the rule-based
and multi-agent pipelines. Upload a screenshot for a media-enabled bug
(see \texttt{data/bug\_with\_media/}).

\subsection{Step 8: (Optional) CLI runs}

Without the UI, both pipelines can be run headless:

\begin{lstlisting}[language=bash]
# Rule-based (fast, deterministic, no API key)
python3 run_agent.py                                # all data/*.txt
python3 run_agent.py data/bug_01_p0_checkout_crash.txt

# Multi-agent (primary, needs Claude CLI + optionally Gemini)
python3 run_multi_agent.py                          # all data/*.txt + bug_with_media/*
python3 run_multi_agent.py data/bug_01_p0_checkout_crash.txt
python3 run_multi_agent.py data/bug_with_media/bug_otp_send_failed
\end{lstlisting}

Both write to \texttt{output/} (rule-based) or \texttt{output\_claude/}
(multi-agent) as \texttt{ticket\_<stem>\_<timestamp>.json}.

\begin{screenshotbox}[Terminal screenshot of a multi-agent run]
Save at \texttt{report/screenshots/06\_cli\_run.png}. Include the full
progress log with all the ``[media]'', ``[classifier]'', ``[similarity]''
and ``[host]'' lines visible.
\end{screenshotbox}

\newpage

% ────────────────────────── SECTION: TESTING ──────────────────────────
\section{Testing and evaluation}

\subsection{Test corpus}

The repository ships with 15 hand-crafted text test bugs under
\texttt{data/*.txt} covering all priority levels and all component labels,
plus 5 media-enabled test bugs under \texttt{data/bug\_with\_media/} (a
mix of images and short videos). Each test bug has:

\begin{itemize}[leftmargin=1.4em,nosep]
\item An expected priority (P0 / P1 / P2)
\item An expected component
\item An expected duplicate (or ``no duplicate'') for a subset of bugs
      that are designed to match known FLIPPI tickets
\end{itemize}

\subsection{Leave-one-out benchmark}

For classifier accuracy I use leave-one-out cross-validation on the full
564-bug corpus. Each bug is held out, the LogReg is retrained on the
other 563, the prediction is compared to the true label. Total accuracy
so far: 69.5\% at threshold 0.50; the 0.60 threshold has not been re-run
through this harness.

\subsection{Regression test folder}

\texttt{tests/} contains 15 paired input/output examples --- one folder
per test with the bug text as \texttt{email.txt}, the rule-based output
as \texttt{triage\_notes.json}, and the multi-agent output as
\texttt{triage\_notes\_claude.json}. A regenerate script
(\texttt{tests/\_build.py} and \texttt{tests/\_run\_claude.py}) rebuilds
everything against the current pipeline. This is not a CI harness --- it's
a visual-diff regression tool for a mentor review.

\newpage

% ────────────────────────── SECTION: LIMITATIONS ──────────────────────────
\section{Limitations and next steps}

Being honest about what's not fixed:

\subsection{Zero training data for the Immersive team}

The 564-bug corpus has no bugs labelled \texttt{immersive}. The classifier
has never seen an immersive bug. The routing table in \texttt{immersive.md}
is derived from the SLAP org chart, not from data. This is the top
data-side gap.

\subsection{No manager mapped for DS}

\texttt{TEAM\_MANAGERS["DS"]} is intentionally unmapped in
\texttt{src/team\_config.py} because the DS manager isn't in the org chart
I have access to. DS bugs where no engineer has similar-bug history return
``no-candidates,'' which the UI surfaces as a manual-triage warning. One
line of config away from being fixed once the name is known.

\subsection{Threshold 0.60 not yet re-benchmarked}

The change from 0.50 $\rightarrow$ 0.60 was based on reasoning, not
measurement. Re-running the LOO harness would give a concrete number to
cite.

\subsection{Backend label cleanup}

The audit found $\sim$60 mis-labelled Backend bugs (should be DS / UI /
Backend-Labs). Fixing those in Jira would jump projected accuracy to
78--82\% with no model changes. This is the single highest-leverage
improvement available.

\subsection{Production gaps (out of scope for the prototype)}

\begin{itemize}[leftmargin=1.4em,nosep]
\item \textbf{Write path}: the ``Publish to Jira'' button is a no-op.
\item \textbf{Authentication}: single-user assumption.
\item \textbf{Monitoring}: no metrics on classifier accuracy in
      production, no alerting on Claude sub-agent failures.
\item \textbf{Scale}: no queuing, no batching, no worker pool.
\item \textbf{Manual JWT refresh}: the Gemini token has to be refreshed
      manually every hour.
\end{itemize}

Each of these is a real production concern but genuinely out of scope for
a prototype whose purpose is to validate the design.

\newpage

% ────────────────────────── SECTION: PROJECT MAPPING ──────────────────────────
\section{Prototype $\rightarrow$ production mapping}

For a mentor asking ``how does this become a real service?'':

\begin{table}[h]
\centering
\small
\begin{tabularx}{\linewidth}{lX}
\toprule
\textbf{Prototype (this repo)} & \textbf{Production (Flipkart PaaS)} \\
\midrule
Bug report \texttt{.txt} file           & Gmail email via \texttt{fk-mart-ai-pulse} \\
Direct Python function calls            & Synapse (auth / routing layer) \\
Python script coordinator               & Astral (agent runtime) \\
\texttt{agent\_parser.py} (rule-based)  & Genvoy / FK-GPT (LLM parsing) \\
\texttt{embedding\_similarity.py}       & Vector One (managed vector DB) \\
\texttt{agent\_scorer.py} (rule-based)  & Genvoy / FK-GPT (LLM scoring) \\
Jira REST API (read)                    & Jira via MART MCP \\
Output JSON file                        & Pulse SMTP reply to reporter \\
\bottomrule
\end{tabularx}
\end{table}

Each prototype component is a one-line-change swap in production ---
the business logic (dedup rules, component routing, ADF structure) stays
the same.

\newpage

% ────────────────────────── APPENDIX ──────────────────────────
\appendix

\section{Sample output JSON}

Output for one of the sample bugs
(\texttt{data/bug\_01\_p0\_checkout\_crash.txt}), truncated for space:

\begin{lstlisting}
{
  "generated_at": "2026-07-08T09:00:00",
  "input_file":   "data/bug_01_p0_checkout_crash.txt",
  "pipeline":     "multi-agent (Astral)",
  "parsed_bug": {
    "title":            "[Checkout]: App crashes on Proceed to Pay",
    "platform":         "Android",
    "app_version":      "2.4.2",
    "component_hint":   "Backend",
    "reproducibility":  "100%",
    "reporter":         "Rahul Verma <rahul.verma@flipkart.com>"
  },
  "jira_ticket_draft": {
    "fields": {
      "project":     {"key": "FLIPPI"},
      "issuetype":   {"id":  "10036"},
      "summary":     "[Checkout]: App crashes on Proceed to Pay ...",
      "priority":    {"id":  "10000"},
      "description": {...ADF...},
      "components":  [{"name": "UI"}],
      "labels":      ["slap", "agentic-triage", "ui", "android"]
    }
  },
  "triage_notes": {
    "team":                    "UI",
    "jira_component":          "UI",
    "component_confidence":    0.71,
    "priority":                "P0",
    "severity_justification":  "100% repro crash on payment button ...",
    "owner_suggestion":        "Hyzam",
    "owner_reason":            "Highest similar-bug match count on UI ...",
    "owner_method":            "claude",
    "duplicate_of":            "FLIPPI-1198",
    "duplicate_confidence":    0.79,
    "similar_bugs":            [ ... top-10 ... ],
    "retrieval": {
      "top_relevance_score":   0.79,
      "is_low_confidence":     false
    },
    "media_findings":          [ ... ],
    "media_combined_summary":  "...",
    "pipeline":                "multi-agent (Astral)",
    "quality_issues":          []
  }
}
\end{lstlisting}

\section{Repository layout}

\begin{lstlisting}
slap-bug-triage/
+-- app.py                      # Streamlit UI
+-- run_multi_agent.py          # Multi-agent CLI runner
+-- run_agent.py                # Rule-based CLI runner
+-- main.py                     # Anthropic SDK pipeline (blocked)
+-- build_embedding_index.py    # One-time index builder
+-- build_repo_skills.py        # Per-repo skill-file generator
+-- CLAUDE.md                   # Project context for future Claude sessions
+-- DOCUMENTATION.md            # Mentor-facing runbook
+-- report/final_report.tex     # This report
+-- src/
|   +-- claude_cli.py           # subprocess wrapper around `claude -p`
|   +-- genvoy_client.py        # Flipkart Gemini proxy wrapper
|   +-- embedding_classifier.py # LogReg + Claude+skills fallback
|   +-- embedding_similarity.py # cosine + cross-encoder rerank
|   +-- team_config.py          # MANAGER_NAMES + TEAM_MANAGERS
|   +-- jira_client.py          # read-only Jira REST wrapper
|   +-- repo_context.py         # repos.json loader
|   +-- agent_ticket_builder.py # ADF ticket builder
|   +-- ... (rule-based path)
|   +-- agents/
|       +-- host_agent.py       # Astral - the coordinator
|       +-- subagent_media.py   # images + videos
|       +-- subagent_parser.py
|       +-- subagent_form_consistency.py
|       +-- subagent_dedup.py
|       +-- subagent_owner.py
|       +-- subagent_triage.py
+-- slap_context/
|   +-- SLAP_KNOWLEDGE.md       # Figma-extracted screen catalog
|   +-- reference_screens/*.png # 16 labelled reference PNGs
|   +-- architecture/
|       +-- repos.json          # team-repo manifest
|       +-- Backend.md, UI.md, ... (team-level skills)
|       +-- repos/
|           +-- edison.md, dropsense.md, ... (per-repo skills)
+-- data/
|   +-- bug_report.txt          # sample inputs
|   +-- bug_01_p0_checkout_crash.txt ... (15 test bugs)
|   +-- bug_with_media/         # 5 test bugs with images / videos
|   +-- embedding_index.npz     # gitignored
|   +-- embedding_index_logreg.pkl  # gitignored
|   +-- embedding_index_team_roster.json  # gitignored
|   +-- corrections.csv         # UI override log; active learning
+-- output/                     # rule-based outputs, gitignored
+-- output_claude/              # multi-agent outputs, gitignored
+-- tests/                      # 15 paired regression tests
+-- .env.example
+-- requirements.txt
\end{lstlisting}

\section{Screenshots and videos --- attachment index}

Place the following files under \texttt{report/screenshots/} and
\texttt{report/videos/}. The document will render them automatically once
the \texttt{includegraphics} lines are uncommented in each section above.

\begin{table}[h]
\centering
\small
\begin{tabular}{ll}
\toprule
\textbf{File} & \textbf{Purpose} \\
\midrule
\texttt{01\_landing.png}         & Streamlit landing page with a bug pasted, no results yet \\
\texttt{02\_pipeline\_running.png} & Live progress stepper mid-run \\
\texttt{03\_media\_findings.png} & Media findings tab for a bug with a screenshot \\
\texttt{04\_ui\_overview.png}    & Full-page UI overview \\
\texttt{05\_owner\_dropdown.png} & Owner dropdown open, managers pinned at the top \\
\texttt{06\_cli\_run.png}        & Terminal screenshot of \texttt{run\_multi\_agent.py} \\
\texttt{07\_token\_banner.png}   & Red GENVOY\_TOKEN-expired banner \\
\texttt{08\_ambiguity\_banner.png} & Ambiguity banner with full probability distribution \\
\midrule
\texttt{videos/demo\_walkthrough.mp4} & 2--3 minute end-to-end demo screen recording \\
\bottomrule
\end{tabular}
\end{table}

\section{Compilation instructions}

To build this report from source:

\begin{lstlisting}[language=bash]
cd report/
pdflatex final_report.tex     # first pass
pdflatex final_report.tex     # second pass for cross-references
# final_report.pdf now exists
\end{lstlisting}

Or with a modern LaTeX distribution:

\begin{lstlisting}[language=bash]
latexmk -pdf final_report.tex
\end{lstlisting}

All packages used are standard (part of MiKTeX / TeX Live) --- no external
tooling required.

\end{document}
